\section{Future  Development Plan}

We plan to continuously update and improve DeXposure-FM. Below is our future development plan.

%\begin{itemize}
    %\item 
    
\subsection{Enlarging Training Data Pool} 
The current DeXposure-FM is trained on weekly inter-protocol snapshots from the DeXposure dataset, while leveraging pretrained TabPFN weights. While this provides broad on-chain coverage at a stable temporal resolution, it omits key off-chain channels through which DeFi is funded, hedged, and linked to traditional markets. In DeXposure-FM~2.0, we plan to expand the training pool by incorporating additional data sources, including centralized exchange (CEX) balances and order-book–implied exposures, over-the-counter (OTC) derivatives and lending agreements, and disclosures on off-chain stablecoin reserves (e.g., USDC’s Treasury-backed holdings). We also aim to add higher-frequency data (daily or intraday, where feasible), particularly around stress episodes, to better capture rapid deleveraging and liquidity dynamics that weekly aggregation smooths out.

\subsection{Expanding Proxies for Credit Exposure}
Total Value Locked (TVL) is a convenient but imperfect proxy for economic exposure. It aggregates heterogeneous assets into a single dollar value without adjusting for liquidity, haircuts, or collateral quality. As a result, equal TVL can imply very different loss-absorbing capacity and contagion risk; for example, \$1B held in thinly traded governance tokens is not economically equivalent to \$1B in USDC, and TVL can also move mechanically with prices even when underlying positions are unchanged. To address these limitations, DeXposure-FM~2.0 will incorporate asset-level risk weights and composition features (e.g., liquidity and volatility proxies, depeg risk, and observable collateral rules) to produce risk-adjusted exposure estimates and more interpretable systemic-risk indicators.

\subsection{Model Drift and Continuous Update}
DeFi evolves rapidly: new chains launch, token standards change, and incentive mechanisms shift. This non-stationarity implies that a model trained on historical regimes may gradually lose accuracy as exposures and flows depart from past patterns, including both gradual structural change and discrete breaks from major upgrades or new collateral rules. Periodic retraining on updated snapshots is therefore essential to maintain forecast quality, ideally paired with drift monitoring (e.g., calibration and error shifts by sector/chain) and versioned releases so that risk signals can be traced to specific data vintages and model iterations.

\subsection{Innovative Model Architecture}
The future DeXposure-FM 2.0 will include an innovative hybrid architecture that combines the current auto-regressive forecasting paradigm \citep{RasulEtAl2023LagLlama} with diffusion-based generative modeling \citep{MeijerChen2024DiffusionTS}. Auto-regressive models are efficient and accurate for point and conditional forecasts, but they can struggle to represent globally coherent network trajectories over long horizons, especially under stress. Diffusion models, by contrast, provide a principled way to learn complex conditional distributions by gradually denoising samples, and have recently shown strong performance in structured generation tasks, such as in diffusion language models \citep{LiEtAl2022DiffusionLM}. %In DeXposure-FM 2.0, we plan to use an autoregressive backbone to produce fast, state-dependent predictions of key sufficient statistics (e.g., protocol states and coarse exposure aggregates), while a conditional diffusion module refines these into full network realizations—jointly sampling link existence and weights in a way that preserves structural constraints such as sparsity, sectoral organization, and budget/flow consistency. This hybrid design would allow the model to output calibrated predictive distributions (not just point forecasts), support scenario-conditional stress testing via controlled guidance, and improve robustness to regime shifts by explicitly modeling the space of plausible future network configurations rather than committing early to a single trajectory.

\subsection{Open Source, Competitions, and Community Development}
We endeavor to develop DeXposure-FM to become fundamental infrastructure rather than a closed research artifact. We have released an open-source version of model weights and training code. We will continue our open-source efforts when improving the model. Based on the DeXposure dataset and DeXposure-FM model, we will organize public benchmarking competitions. These challenges will be designed to encourage robust methods and will include baseline implementations, leaderboards, and model cards reporting calibration and failure modes. We also aim to build a sustained community around DeFi economic measurement. By lowering barriers to entry and making the measurement process transparent, we hope to accelerate cumulative progress on reliable, policy-relevant tools for monitoring systemic risk in decentralized financial networks.



\section{Conclusions}
\label{sec:conclusion}

Credit exposures in decentralized finance (DeFi) are often embedded in token holdings and smart-contract interactions, rather than recorded as explicit bilateral claims, forming a tightly coupled network of inter-protocol dependencies. In such an environment, a disturbance to a single token can propagate across protocols, generating sizable and difficult-to-contain contagion. As DeFi becomes more intertwined with traditional financial infrastructure, most prominently via stablecoins, this evolving risk landscape calls for stronger tools to measure and forecast exposures. We introduce DeXposure-FM, which is arguably the first time-series, graph foundation model, designed for DeFi. DeXposure-FM combines a graph–tabular encoder with pretrained weights as initialization and multiple task-specific heads. It is trained on the DeXposure dataset, comprising 43.7 million observations spanning 4,300+ protocols on 602 blockchains and 24,300+ unique tokens. The model is trained for credit-exposure forecasting, jointly learning the dynamics of protocol-level flows, the evolving topology and weights of exposure links, and sectoral structure (e.g., lending, exchanges, asset management, bridges, stablecoins). Across a suite of machine-learning benchmarks, DeXposure-FM consistently outperforms strong competitors, including a prior graph foundation model and state-of-the-art temporal graph neural networks. We also translate the model outputs into economically interpretable indicators, including time-varying measures of systemic importance, sector-to-sector spillover indices, and early-warning signals based on shifts in network concentration and reliance on fragile collateral, which illustrate applications to macroprudential monitoring, DeFi stress testing, and counterfactual policy analysis. The tools are fully supported in he empirical validations. The model and code have been publicly available.


%DeXposure-FM demonstrates that graph foundation models can effectively measure and forecast inter-protocol credit exposures in decentralised finance. Trained on the DeXposure dataset spanning over 4,300 protocols across 602 blockchains from 2020 to 2025, the model learns joint representations of protocol-level dynamics, edge-level exposures, and sector structure. Evaluated on multi-step forecasting, historical stress-event analysis, and missing-edge imputation, GraphPFN-based variants maintain stable predictive performance (AUROC $>$0.91) across forecast horizons up to 14 weeks, outperforming temporal GNN baselines while prioritising economically significant edges.

%These predictive capabilities enable broader applications. DeXposure-FM provides a measurement layer for macroprudential monitoring: systemic importance scores, sector spillover matrices, and early-warning indicators derived from network structure. These outputs can support regulators in identifying concentration risks and conducting stress tests on DeFi infrastructure.

%While the current model is limited to on-chain DeFi data at weekly resolution, these constraints point toward promising extensions: integrating off-chain data sources to bridge DeFi and traditional finance, incorporating asset-level risk weights for finer exposure quantification, and scaling to real-time monitoring. As DeFi infrastructure continues to grow and interconnect with the broader financial system, we believe graph-based foundation models like DeXposure-FM can serve as a critical building block for cross-ecosystem risk surveillance.
